\subsubsubsection{Mechanical}
According to the Mechanical Department, the current payload bay (as of March 2026) can house a full-width board (~87x87 mm) that is approximately 35 mm tall.
Moreover, the camera can be mounted on an angle to minimise its height if the lens is too large.

\subsubsubsection{Communications}
As of March 2026, the Comms department estimates 7 to 18 minutes of contact with the CubeSat for a single ground station, averaged to 10 minutes, with two passes per day.
They estimated a reliable download rate of 3 kb/s, with optimal speeds of 6 kb/s, and a potential maximum of 40 kb/s.
Through a verbal agreement with Mission, we concluded that a three day window would be the longest time for transmitting a single image, with a single day window being an ideal time to transmit an image.
Moreover, during this window, approximately 60\% of the data would be dedicated to the photo, where the other 40\% is for telemetry data.
A breakdown of the estimated download sizes for a single ground station can be seen in \autoref{tab:v1-comms-reqs}.

\begin{indentedfigure}
    \centering
    \begin{tabular}{|c|c|c|}
        \hline
        \textbf{Rate} & \textbf{1 Day} & \textbf{3 Days} \\
        \hline
        3 kb/s & 270 kB & 810 KB \\
        \hline
        6 kb/s & 540 kB & 1.62 MB \\
        \hline
        40 kb/s & 3.6 MB & 10.8 MB \\
        \hline
    \end{tabular}
    \caption{Estimated download sizes for a single ground station at 60\% usage}
    \label{tab:v1-comms-reqs}
\end{indentedfigure}

While the final compression algorithm is pending, initial estimates using \cite{balalam_image_compression} and \cite{squoosh_editor} suggest a 10:1 JPEG compression ratio is achievable before significant artifacting occurs.
Given the operational similarities between JPEG and ICER, this figure serves as a reliable baseline for the current design phase.
Assuming the raw data is received in the YCbCr format at 12-bit resolution (see \autoref{sec:appendices} for more on YCbCr), the download sizes in \autoref{tab:v1-comms-reqs} can be converted into approximate CIS pixel counts.

\begin{indentedfigure}
    \centering
    \begin{tabular}{|c|c|c|}
        \hline
        \textbf{Rate} & \textbf{1 Day} & \textbf{3 Days} \\
        \hline
        3 kb/s & 1.8 MP & 5.4 MP \\
        \hline
        6 kb/s & 3.6 MP & 10.8 MP \\
        \hline
        40 kb/s & 24 MP & 72 MP \\
        \hline
    \end{tabular}
    \caption{Approximate megapixel count for assumed Comms requirements}
    \label{tab:v1-pixel-count}
\end{indentedfigure}

Based on these benchmarks, a CIS under 5 MP optimises the trade-off between visual fidelity and transmission overhead.
This resolution supports QuadHD (2K) images with flexibility for 2x binning to HD.
While higher data rates or multiple ground stations would allow for lower compression ratios, the pixel count remains viable under degraded conditions, capable of downlinking a 5 MP frame over 72 hours at a 3 kb/s baseline.
This approach helps prioritise mission reliability over ideal-case performance.